\textbf{Phần 1. Trắc nghiệm lựa chọn}
\setcounterpageref{ex}{0}

\begin{ex} %[3D6N36]
Kết quả của phép tính $100 \times 4$ là
\choice
{$300$}
{$500$}
{$450$}
{\True $400$}
\loigiai{
$100 \times 4 = 400$.\\
\textbf{Đáp án của câu hỏi là D}
}
\end{ex}

\begin{ex} %[3D6N36]
Kết quả của phép nhân không nhớ $123 \times 3$ là
\choice
{\True $369$}
{$359$}
{$379$}
{$349$}
\loigiai{
$123 \times 3 = 369$.\\
\textbf{Đáp án của câu hỏi là A}
}
\end{ex}

\begin{ex} %[3D6H36]
Kết quả của phép nhân có nhớ $142 \times 4$ là
\choice
{$558$}
{$578$}
{\True $568$}
{$468$}
\loigiai{
$142 \times 4 = 568$.\\
\textbf{Đáp án của câu hỏi là C}
}
\end{ex}

\begin{ex} %[3D6H36]
Kết quả của phép nhân $215 \times 3$ là
\choice
{$635$}
{$655$}
{$615$}
{\True $645$}
\loigiai{
$215 \times 3 = 645$.\\
\textbf{Đáp án của câu hỏi là D}
}
\end{ex}

\begin{ex} %[3D6H36]
Tìm $x$ biết: $x : 3 = 152$.
\choice
{$446$}
{\True $456$}
{$466$}
{$450$}
\loigiai{
$x = 152 \times 3 = 456$.\\
\textbf{Đáp án của câu hỏi là B}
}
\end{ex}

\begin{ex} %[3D6H36]
Tính giá trị biểu thức: $120 \times 4 + 150$.
\choice
{\True $630$}
{$610$}
{$650$}
{$600$}
\loigiai{
$120 \times 4 + 150 = 480 + 150 = 630$.\\
\textbf{Đáp án của câu hỏi là A}
}
\end{ex}

\begin{ex} %[3D6V36]
Một trường tiểu học có 3 lớp Ba, mỗi lớp có 32 học sinh. Tổng số học sinh lớp Ba của trường đó là
\choice
{$92$ học sinh}
{\True $96$ học sinh}
{$98$ học sinh}
{$102$ học sinh}
\loigiai{
$32 \times 3 = 96$ học sinh.\\
\textbf{Đáp án của câu hỏi là B}
}
\end{ex}

\begin{ex} %[3D6V36]
Mỗi bao gạo nặng $125\text{ kg}$. Hỏi 4 bao gạo như thế nặng tất cả bao nhiêu ki-lô-gam?
\choice
{$450\text{ kg}$}
{$550\text{ kg}$}
{\True $500\text{ kg}$}
{$600\text{ kg}$}
\loigiai{
$125 \times 4 = 500\text{ kg}$.\\
\textbf{Đáp án của câu hỏi là C}
}
\end{ex}

\begin{ex} %[3D6V36]
Một cuộn dây điện dài $230\text{ m}$. Hỏi 3 cuộn dây điện như thế dài bao nhiêu mét?
\choice
{$660\text{ m}$}
{$720\text{ m}$}
{$750\text{ m}$}
{\True $690\text{ m}$}
\loigiai{
$230 \times 3 = 690\text{ m}$.\\
\textbf{Đáp án của câu hỏi là D}
}
\end{ex}

\begin{ex} %[3D6V36]
Một hình vuông có cạnh dài $145\text{ cm}$. Chu vi hình vuông đó là
\choice
{\True $580\text{ cm}$}
{$560\text{ cm}$}
{$600\text{ cm}$}
{$540\text{ cm}$}
\loigiai{
$145 \times 4 = 580\text{ cm}$.\\
\textbf{Đáp án của câu hỏi là A}
}
\end{ex}

\begin{ex} %[3D6V36]
Tính giá trị biểu thức: $118 \times 5 - 190$.
\choice
{$380$}
{\True $400$}
{$420$}
{$450$}
\loigiai{
$118 \times 5 - 190 = 590 - 190 = 400$.\\
\textbf{Đáp án của câu hỏi là B}
}
\end{ex}

\begin{ex} %[3D6C36]
Một cửa hàng nhập về 3 thùng dầu, mỗi thùng chứa $150\text{ l}$. Cửa hàng đã bán được $280\text{ l}$. Lượng dầu còn lại là
\choice
{$160\text{ l}$}
{$180\text{ l}$}
{\True $170\text{ l}$}
{$190\text{ l}$}
\loigiai{
Tổng số dầu: $150 \times 3 = 450\text{ l}$.\\
Còn lại: $450 - 280 = 170\text{ l}$.\\
\textbf{Đáp án của câu hỏi là C}
}
\end{ex}

\textbf{Phần 2. Trắc nghiệm đúng sai}
\setcounterpageref{ex}{0}

\begin{ex} %[3D6H36]
Xét tính đúng sai của các phép nhân:
\choiceTF
{\True $112 \times 4 = 448$.}
{\True $235 \times 2 = 470$.}
{$108 \times 5 = 530$.}
{\True $108 \times 5 = 540$.}
\loigiai{
\begin{itemize}
    \item $112 \times 4 = 448$: Đúng.
    \item $235 \times 2 = 470$: Đúng.
    \item $108 \times 5 = 540 \neq 530$. Ý 3 Sai, Ý 4 Đúng.
\end{itemize}
\begin{itemize}
    \item a) \textbf{Đúng}.
    \item b) \textbf{Đúng}.
    \item c) \textbf{Sai}.
    \item d) \textbf{Đúng}.
\end{itemize}
}
\end{ex}

\begin{ex} %[3D6V36]
Một trang trại nuôi 3 chuồng vịt, mỗi chuồng có 160 con vịt. Xét tính đúng sai:
\choiceTF
{\True Cả 3 chuồng nuôi tất cả 480 con vịt.}
{\True Nếu bán đi 150 con vịt thì còn lại 330 con vịt.}
{Nếu nhập thêm 50 con vịt thì tổng số vịt là 520 con.}
{\True Nếu nhập thêm 50 con vịt thì tổng số vịt là 530 con.}
\loigiai{
\begin{itemize}
    \item Tổng: $160 \times 3 = 480$ con. Đúng.
    \item Còn lại: $480 - 150 = 330$ con. Đúng.
    \item Thêm 50: $480 + 50 = 530$ con $\neq 520$. Ý 3 Sai, Ý 4 Đúng.
\end{itemize}
\begin{itemize}
    \item a) \textbf{Đúng}.
    \item b) \textbf{Đúng}.
    \item c) \textbf{Sai}.
    \item d) \textbf{Đúng}.
\end{itemize}
}
\end{ex}

\textbf{Phần 3. Trả lời ngắn}
\setcounterpageref{ex}{0}

\begin{ex} %[3D6H36]
Tính tích: $132 \times 3$.
\shortans{396}
\loigiai{
\textbf{Đáp số:} $396$\\
$132 \times 3 = 396$.
}
\end{ex}

\begin{ex} %[3D6H36]
Tính tích: $214 \times 4$.
\shortans{856}
\loigiai{
\textbf{Đáp số:} $856$\\
$214 \times 4 = 856$.
}
\end{ex}

\begin{ex} %[3D6V36]
Tính giá trị biểu thức: $125 \times 4 - 200$.
\shortans{300}
\loigiai{
\textbf{Đáp số:} $300$\\
$500 - 200 = 300$.
}
\end{ex}

\begin{ex} %[3D6V36]
Một rạp chiếu phim có 4 hàng ghế, mỗi hàng có 115 chiếc ghế. Tính tổng số ghế trong rạp.
\shortans{460}
\loigiai{
\textbf{Đáp số:} $460$\\
$115 \times 4 = 460$ chiếc ghế.
}
\end{ex}

\begin{ex} %[3D6V36]
Tìm $x$ biết: $x : 5 = 140$.
\shortans{700}
\loigiai{
\textbf{Đáp số:} $700$\\
$x = 140 \times 5 = 700$.
}
\end{ex}

\begin{ex} %[3D6C36]
Một hình chữ nhật có chiều rộng $112\text{ m}$, chiều dài gấp 3 lần chiều rộng. Tính chiều dài hình chữ nhật theo mét.
\shortans{336}
\loigiai{
\textbf{Đáp số:} $336$\\
$112 \times 3 = 336\text{ m}$.
}
\end{ex}
